\documentclass[11pt]{article}
\usepackage[utf8]{inputenc}
\usepackage{amsmath,amssymb,amsthm}
\usepackage[margin=1in]{geometry}
\usepackage{hyperref}
\usepackage{xcolor}
\usepackage{enumitem}
\usepackage{newunicodechar}
\usepackage{booktabs}
\newunicodechar{ℝ}{\ensuremath{\mathbb{R}}}
\newunicodechar{⟨}{\ensuremath{\langle}}
\newunicodechar{⟩}{\ensuremath{\rangle}}
\newunicodechar{ε}{\ensuremath{\varepsilon}}
\newunicodechar{λ}{\ensuremath{\lambda}}

\newcommand{\userbox}[1]{\par\medskip\begin{quote}\small\textbf{\textcolor{blue!60}{DM:}} #1\end{quote}\medskip}
\newcommand{\claudebox}[1]{\par\medskip\begin{quote}\small\textbf{\textcolor{orange!60!black}{Claude:}} #1\end{quote}\medskip}
\newcommand{\gptbox}[1]{\par\medskip\begin{quote}\small\textbf{\textcolor{green!50!black}{GPT-5.5 Pro:}} #1\end{quote}\medskip}

\title{Tide-loop retrospective:\\five formalisations, a side excursion,\\
and the first cross-repo step}
\author{Daniel Murfet}
\date{1 May 2026}

\begin{document}
\maketitle

\begin{abstract}
This document is a retrospective on five Tide-loop excursions and one
adjacent numerical/theoretical exploration carried out across two
afternoons (1--2 May 2026), building on the morning recorded in
\texttt{retrospective.tex} and the few-day saga in
\texttt{retrospective2.tex}. The Tide loop is a thin, deliberation-first
procedure for autonomous formalisation: a human-provided direction;
a Claude--GPT-5.5~Pro vote on a minimal candidate lemma; hand-off to the
\texttt{lean-formalisation} skill until success. Three loops on day~1
produced 1099 lines of new Lean across three branches of the
\texttt{timaeus-research/laplace} repository, plus an adjacent
numerical-and-theoretical exploration with a Python script, three
figures, and three new appendix sections in the susceptibility primer.
A fourth loop the same evening added the rate-aware large-$T$
asymptotic of the universal scaling function. A fifth loop on day~2,
the first to instantiate the framework on a Lean repo other than
\texttt{laplace}, formalised the vanishing of the third cumulant under
the harmonic Gibbs measure in \texttt{timaeus-research/threepoint}.
All five excursions zero \texttt{sorry} / \texttt{axiom} /
\texttt{native\_decide}; total new Lean is 1414 lines across four
branches in \texttt{laplace} (now merged) plus 126 lines in
\texttt{threepoint} (now merged). The aim of this retrospective is to
record what worked about the loop, where it was stress-tested by
repeated use, what skill iteration emerged, and what we learned by
running the loop on a second repo.
\end{abstract}

\tableofcontents

\section{Setting}

The morning of 27 April, recorded in \texttt{retrospective.tex},
formalised the primer's anharmonic covariance formula in 3279 lines of
Lean. The few-day saga that followed, recorded in
\texttt{retrospective2.tex}, brought the multivariate sharp covariance
to the same bar. The framework articulated in DM's notes
``Formalisation First Steps'' (27/28 April) and ``Formalisation Second
Steps'' (28 April) projected this work forward into a recurring loop:

\begin{quote}
\itshape
I'd like to wake up to a new paper with a link in the margin to the
formalisations.
\end{quote}

The present document records the first afternoon during which that loop
ran end-to-end. By the start of the afternoon DM had drafted a small
SRI skill, \texttt{tide}, that codifies the procedure in three steps:

\begin{enumerate}
\item Record the human direction in a tide log entry.
\item Deliberate with GPT-5.5~Pro until both models vote for the same
minimal candidate lemma.
\item Hand off to the \texttt{lean-formalisation} skill until the
proofs land.
\end{enumerate}

Three loops were run, each starting from the previous one's commit
(which is the central image of the loop: a tide of agents spreading
outward from already-formalised territory, with the shoreline rising
each time).\footnote{The Tide-loop framing carries a small vocabulary
of metaphors. The \emph{seabed} is the existing formalised material
that a Tide step launches from: at the start of the afternoon it was
the laplace repository's \texttt{main} branch; after Tide 1 it
included \texttt{Quartic.lean} as well; after Tide 2 it included the
2D semi-degenerate file; and so on. The \emph{shoreline} is the
just-beyond-the-seabed watermark of theorems that are reachable in
one Tide step from the current seabed; each successful step extends
it. The \emph{tide} is the agent activity itself, sweeping over the
seabed and (when systematisation lands the new pieces in
\texttt{lean-common}) genuinely raising the floor for the next
agent to launch from. This vocabulary is from the automation project
README; it is loose but useful.}
The rest of this document records what the runs looked
like, how the deliberation step calibrated, where the skill was
stress-tested, and what was learned.

\section{The five excursions}\label{sec:excursions}

\begin{table}[h]
\centering
\renewcommand{\arraystretch}{1.2}
\begin{tabular}{l l l r r}
\toprule
\textbf{Tide step} & \textbf{Repo} & \textbf{Direction} & \textbf{Lines} & \textbf{Build}\\
\midrule
1. \texttt{quartic-moments}    & \texttt{laplace}    & ``extend to simple degenerate cases'' & 379 & 4s\\
2. \texttt{2d-semi-degenerate} & \texttt{laplace}    & ``Next''  & 514 & 12s\\
3. \texttt{universal-scaling}  & \texttt{laplace}    & ``formalise the universal scaling formula'' & 206 & 8s\\
4. \texttt{f1-large-t}         & \texttt{laplace}    & ``proceed with (a)'' (large-$T$ asymptotic) & 189 & 8s\\
5. \texttt{kappa3-harmonic}    & \texttt{threepoint} & ``$\kappa_3$ vanishes for harmonic Gibbs'' & 126 & 31s\\
\midrule
\multicolumn{3}{r}{\emph{Total new Lean for tide loop}} & \textbf{1414} & \\
\bottomrule
\end{tabular}
\caption{The five Tide-loop excursions of 1--2 May 2026. The Tide-step
identifier doubles as the branch name (\texttt{tide/<id>}). Branches
1--3 are linearly rebased on each other within \texttt{laplace}; branch
4 attaches to \texttt{laplace/main} after 1--3 merged; branch 5 is the
first cross-repo step, on \texttt{threepoint}.}
\label{tab:excursions}
\end{table}

The targets, in plain English:\footnote{Dan: lol}

\begin{itemize}
\item \textbf{Tide 1 --- pure-quartic Gibbs moments.} For
$L(x) = x^4/24$ (genuinely degenerate at the origin: $L''(0) = 0$),
prove the closed forms $\mathbb{E}_t[X^{2n}] = (24/t)^{n/2}\,
\Gamma((2n+1)/4)/\Gamma(1/4)$ for the moment family, and as a
corollary the affine covariance $\operatorname{Cov}_t[ax+c, bx+d]
= ab\,\sqrt{24/t}\,\Gamma(3/4)/\Gamma(1/4)$. The headline rate is
$t^{-1/2}$, contrasting with the $t^{-1}$ of the seabed's nondegenerate
Hessian case.
\item \textbf{Tide 2 --- 2D semi-degenerate covariance.} For
$L(x,y) = (\lambda/2)y^2 + x^4/24$, prove the exact affine covariance
$a_1 b_1 \sqrt{24/t}\,\Gamma(3/4)/\Gamma(1/4) + a_2 b_2/(\lambda t)$,
combining the $t^{-1/2}$ rate of the degenerate quartic direction with
the $t^{-1}$ rate of the nondegenerate quadratic direction.
\item \textbf{Tide 3 --- universal scaling identity.} For
$L^{\rm nd}_\varepsilon(x) = (\varepsilon/2)x^2 + x^4/24$, prove the
exact substitution identity $\langle x^{2n}\rangle_{\mu_t} =
\varepsilon^n \cdot \langle y^{2n}\rangle^V_{\varepsilon^2 t}$ where
$V(y) = y^2/2 + y^4/24$ is the universal potential. This is the
formalisation of the pre-asymptotic crossover phenomenon explored
numerically in the side excursion (\S\ref{sec:near-degenerate}).
\end{itemize}

Each step's seabed is the previous step's tip. The branch architecture
is therefore linear: \texttt{tide/quartic-moments} branches off
\texttt{main}; \texttt{tide/2d-semi-degenerate} branches off
\texttt{tide/quartic-moments}; \texttt{tide/universal-scaling} branches
off \texttt{tide/2d-semi-degenerate}. The primer pins all three sets of
margin markers to a single \texttt{\textbackslash quarticCommit} (the
tip of the most recent branch), which collapses to
\texttt{\textbackslash laplaceCommit} once everything merges.

\section{Anatomy of a Tide step}

The skill makes the procedure explicit. We illustrate it with Tide 1
(\texttt{quartic-moments}); Tide 2 and 3 followed the same shape.

\subsection{Step 1 --- direction}

DM's direction:

\userbox{The primer work in the laplace repository formalises some
statements about covariances against a posterior where $L$ is
nondegenerate. It would be useful to do this also in some simple
degenerate cases.}

The skill prescribes opening a tide log file\footnote{For Tide 1, this
is \texttt{sri/projects/primer/tide-log/2026-05-01-tide-degenerate-quartic-cov.md}.}
with the direction recorded verbatim, the seabed pointer
(\texttt{laplace}, commit SHA at start), and the start date.

\subsection{Step 2 --- deliberation}

\paragraph{Drafting.} Claude reads the seabed (theorem signatures and
\texttt{CLAUDE.md} only, not full proofs) and drafts 2--3 candidates.
For Tide 1 the candidates were:

\begin{itemize}
\item A: a closed-form identity for $\mathbb{E}_t[X^2]$ under the pure
quartic, valid for all $t > 0$;
\item B: an asymptotic for $\operatorname{Cov}_t[\phi,\psi]$ under the
pure quartic with general $C^2$ observables;
\item C: a 2D semi-degenerate covariance combining quartic and
quadratic.
\end{itemize}

\paragraph{Consulting.} Claude sends all three candidates to GPT-5.5
Pro via the \texttt{timaeus-research} skill's \texttt{query\_gpt.py},
asking three explicit questions: are the statements correct, which is
the minimal good target, are there missed candidates? GPT's response is
saved verbatim to the log under \texttt{\#\# GPT-5.5~Pro v1}.

For Tide 1, GPT delivered three useful corrections:

\begin{enumerate}
\item Candidate B was \emph{wrong as written}: with only a $C^2$
remainder bound the cross-terms produce a parasitic $O(t^{-3/4})$
contribution rather than $O(t^{-1})$; recovering the cleaner remainder
needs $C^3$ smoothness so that cubic terms cancel by parity. (This
mirrors the regularity-staircase familiar from
\texttt{lem:laplace\_cov2} in the seabed.)

\item Candidate C had \emph{hidden coefficients}: the writeup gave an
explicit $a_2 b_2/(\lambda t)$ term, but with only an $O(\|(x,y)\|^2)$
remainder hypothesis on the observables, additional $t^{-1}$
contributions from the quartic direction's higher jets are absorbed
into the remainder, making the explicit coefficient uninterpretable
without strengthening the hypothesis.

\item DM's offhand mention of $L = x^2 y^4$ as a candidate potential
was \emph{not even integrable}: $e^{-t x^2 y^4}$ does not define a
finite Gibbs posterior on $\mathbb{R}^2$, since on the axes the
exponent is identically zero.
\end{enumerate}

In addition GPT proposed a strictly stronger version of A (call it
A$''$): the full even-moment family $\mathbb{E}_t[X^{2n}]$ in closed
form, at no extra proof cost over A, providing a reusable
moment API. Claude adopted this in v2.

\paragraph{Voting.} The skill mandates an explicit two-way vote. For
Tide 1, both Claude and GPT-5.5~Pro voted for A$''$ + an affine
covariance corollary. Vote recorded in the log:

\begin{quote}
\#\# Vote
\par
- Claude: A$''$ + A$'$ (moment family + affine covariance corollary)\par
- GPT-5.5~Pro: A$''$ (with A$'$ as the closely-related corollary).
\end{quote}

\paragraph{Numerical sanity.} Before handing off to Step 3, Claude runs
a numerical sanity check using \texttt{scipy.integrate.quad} to confirm
the closed-form identity to machine precision at several parameter
values. For Tide 1: 24 cases ($n = 0,\ldots,6$, $t \in \{0.5, 1, 5, 100\}$),
all matching to $|\text{diff}| < 10^{-13}$.

\subsection{Step 3 --- formalisation}

The agreed candidate is handed off to the \texttt{lean-formalisation}
skill. Tide 1's hand-off proposed file path
(\texttt{Laplace/OneD/Quartic.lean}), theorem names
(\texttt{quartic\_partition}, \texttt{quartic\_moment\_even},
\texttt{quartic\_cov\_affine}, etc.), proof skeleton (substitution
$u = (t/24)^{1/4} x$ to Mathlib's
\texttt{integral\_rpow\_mul\_exp\_neg\_mul\_rpow}, then evenness
doubling for the full line), and the closed-form expectation from the
numerical check.

For Tide 1 the skeleton typechecked first try, and the proofs filled
in by direct application of the existing seabed pattern (the file
\texttt{Laplace/OneD/GaussianMoments.lean}, which formalises the
analogous result for $L = x^2/2$). The only non-routine piece was
proving \texttt{quartic\_integrable\_pow}: full-line integrability of
$x^n e^{-t x^4 / 24}$. This was solved by Gaussian comparison via the
square-completion identity
\[
\frac{t x^4}{24} - x^2 = \frac{(t x^2 - 12)^2}{24 t} - \frac{6}{t}
\ge -\frac{6}{t},
\]
which dominates the integrand by a polynomial-times-Gaussian and
brings it under Mathlib's
\texttt{integrable\_rpow\_mul\_exp\_neg\_mul\_sq}. A second GPT consult
was started for this lemma in parallel with Claude's local attempt;
Claude's local proof landed before GPT's response (which timed out at
the 15-minute API limit), confirming the skill's advice that for
narrowly-scoped tactical Lean questions one should race rather than
block on GPT.

\section{The deliberation step in practice}

Tide 2 and Tide 3 followed the same template, with patterns emerging
across the three rounds.

\paragraph{GPT consistently delivered a strict improvement.} In each
round, GPT-5.5 Pro proposed a generalisation of Claude's drafted
candidate that strictly subsumed the original at zero extra proof
cost.

\begin{itemize}
\item Tide 1: A (closed-form $\mathbb{E}_t[X^2]$) $\to$ A$''$ (full
moment family).
\item Tide 2: A (2D semi-degenerate affine covariance) was already the
right headline, but GPT recommended an internal mixed-monomial
factorisation lemma $\langle z_1^m z_2^n\rangle = \langle x^m\rangle
\langle y^n\rangle$ as the load-bearing helper, instead of the six
specialised cases Claude had drafted. The lemma is a strict
generalisation costing one proof line beyond the specific cases.
\item Tide 3: C (substitution lemma for the monomial integrand) $\to$ D
(substitution lemma for arbitrary $f$), again a strict generalisation
at zero extra proof cost.
\end{itemize}

The pattern is consistent enough to be a Tide-skill rule: \emph{when
drafting a candidate, draft the most specific true thing; expect GPT to
broaden it; if GPT does not, draft B or C as the broader version
yourself}. The deliberation is reliably one round (Claude $\to$ GPT
$\to$ vote), not the four rounds the skill caps at.

\paragraph{GPT's bug-catching is asymmetrically valuable.} Tide 1's
correction of Candidate B (the spurious $O(t^{-1})$ remainder) and the
flag of $L = x^2 y^4$ as non-integrable were the most valuable
contributions of the entire afternoon. Both would have wasted hours of
Lean work if not caught at the deliberation stage. The
\emph{value} of Step 2 is precisely that bugs land in deliberation
words (cheap) rather than in proof state (expensive).

\paragraph{The vote is genuinely a vote, not a rubber stamp.} Tide 2's
deliberation produced two materially different positions (Claude
proposed thin 2D wrappers using \texttt{volume.prod}; GPT proposed
iterated integrals). Both were valid; the vote was on the headline
target, with the wrapper choice noted as ``Claude diverges from GPT's
recommendation for $\sim 1$ implementation reason'' rather than a
re-deliberation. The skill's rule (``A round ends with both parties
stating which single candidate they back'') is loose enough to absorb
this kind of micro-divergence without re-running.

\paragraph{Rounds 2--4 of the skill were never used.} Across three
Tides, every deliberation converged in one round. A second round may
become useful when candidates are at a level of mathematical depth
where GPT's first-pass corrections are not exhaustive; for the present
class of substitution-and-Gaussian-moment results, one round is the
norm.

\section{The formalisation step in practice}

Across the three Tides, the formalisation step was uneven in difficulty,
broadly in proportion to the size of the step outward.

\paragraph{Tide 1: routine, with one minor escalation.} Most of
\texttt{Quartic.lean} mirrors \texttt{GaussianMoments.lean} from the
seabed. The only nontrivial piece, full-line integrability, prompted a
GPT consult that timed out before Claude finished a local proof. Total
attention: $\sim 1$ hour.

\paragraph{Tide 2: longer, but no GPT consults.} The 2D step required
new infrastructure: thin 2D Gibbs wrappers, partition-function
factorisation, mixed-monomial factorisation, six 2D atom
integrabilities (built from \texttt{Integrable.mul\_prod}), and a
chained \texttt{integral\_add} expansion for the affine covariance.
The known Pi.add gotcha from the seabed's \texttt{CLAUDE.md} bit once
and was resolved by the documented type-ascribed-single-lambda fix.
Total attention: $\sim 1.5$ hours.

\paragraph{Tide 3: short, no GPT consults.} The substitution proof
landed in 30 minutes. The architecture GPT recommended in deliberation
(thin abbreviations + algebraic identity + general substitution lemma
+ partition-function corollary + Gibbs-expectation transport +
headline-at-monomial) collapsed to 206 lines, of which the four
named theorems took $\sim 50$ lines each.

\paragraph{Pattern: zero GPT consults during proof for Tides 2 and 3.}
The deliberation step preceding the formalisation produced enough
strategic alignment that no further consults were needed. This differs
from \texttt{retrospective.tex}'s pattern (three load-bearing GPT
consultations during the formalisation itself), and the difference is
diagnostic: when the Tide deliberation has done its job, the
formalisation is mechanical.

\paragraph{Skill iteration during the run.} The
\texttt{lean-formalisation} skill, authored before the morning's work,
was iterated twice during the afternoon as patterns emerged:
\begin{itemize}
\item After Tide 1, the integrability-via-square-completion pattern
was a natural addition to the skill's ``Patterns and pitfalls''
section (alongside the existing entries on over-engineering, Pi.add,
and template-vs-custom proofs).
\item After Tide 2, the ``\texttt{div\_eq\_zero\_iff} + resolve\_right
to extract numerator-zero from expectation-zero'' move (used three
times in the affine covariance proof) was promoted from a one-off to
a documented idiom.
\end{itemize}
These additions were small (a few sentences each) but represent the
mechanism by which the skill compounds. The
``rising tide'' image is exact: the skill is the seabed; each tide
step adds a thin layer.

\section{Cross-step compounding}

The most striking feature of running three Tide steps in a row is how
much each step uses material from the previous one. Tide 2 imports
\texttt{Laplace/OneD/Quartic.lean} (Tide 1) and reuses
\texttt{quartic\_partition}, \texttt{quartic\_moment\_even},
\texttt{quartic\_expected\_value\_sq}, and
\texttt{quartic\_integrable\_pow\_pot}. Tide 3 imports
\texttt{Laplace/OneD/Anharmonic.lean} (already in the seabed) but
restructures \texttt{nearDegeneratePotential = anharmonicPotential
\_ 0 1} as an abbreviation, immediately giving access to the existing
algebraic and integrability results. The cumulative effect is that
each step is shorter than its predecessor would have been if attempted
from scratch:

\begin{table}[h]
\centering
\renewcommand{\arraystretch}{1.2}
\begin{tabular}{l r r r}
\toprule
& \textbf{New lines} & \textbf{Lines reused} & \textbf{Cumulative library}\\
\midrule
Seabed (start)        &     0 & --      & $\sim$15800\\
Tide 1                &   379 & seabed  &  16179\\
Tide 2                &   514 & + Tide 1 &  16693\\
Tide 3                &   206 & + Tide 1, 2 & 16899\\
\bottomrule
\end{tabular}
\caption{Each step is shorter than the previous in absolute terms, while
the cumulative library grows steadily. Tide 3 in particular is small
because the seabed's anharmonic infrastructure does most of the work.}
\label{tab:lines}
\end{table}

This is what ``the tide rises'' looks like in practice: the seabed
gets thicker, and the next step launches from higher ground. The
\emph{shoreline} (the theorems-just-beyond-reach watermark) advances
visibly. After Tide 1, the shoreline includes 2D semi-degenerate cases.
After Tide 2, the shoreline includes the universal scaling identity.
After Tide 3, the shoreline includes the two-sided asymptote of
$f_n(T)$ (Candidate (2) from the near-degenerate note) and a
crossover-aware uniform-in-$\varepsilon$ inequality (Candidate (3)).
Every successful step makes the next step's deliberation cheaper.

\section{The near-degenerate side excursion}\label{sec:near-degenerate}

Between Tide 2 and Tide 3, DM asked an exploratory question rather
than directing a Tide step:

\userbox{Here's a question: when $L^{\rm nd}$ is nondegenerate but is
``close'' to a degenerate potential $L$ it can still be the case that
asymptotic behaviour of expectation values for $L^{\rm nd}$ are
governed by those for $L$ until $t$ is large. Explore this effect
numerically and theoretically.}

The investigation took $\sim 30$ minutes. The
result\footnote{Recorded at
\texttt{sri/projects/primer/tide-log/2026-05-01-near-degenerate-crossover.md}.}
is the cleanest scaling-collapse picture of the whole afternoon:
under the rescaling $x = \sqrt{\varepsilon} y$ with $T = \varepsilon^2
t$, the moment data for $L^{\rm nd}_\varepsilon(x) =
(\varepsilon/2) x^2 + x^4/24$ collapses onto a universal scaling
function $f_n(T)$ for the fixed potential $V(y) = y^2/2 + y^4/24$.
Below the crossover scale $t_c \sim 1/\varepsilon^2$ the moments
track the pure-quartic asymptotics; above, they track the standard
Laplace asymptotics; the crossover region is the order-of-magnitude
window around $T = 1$.

This was \emph{not} a Tide step. It was a numerical exploration with a
theoretical scaling argument written up alongside the figures, ending
with three explicit ``possible Tide-step formalisations'' of which
candidate (1) became Tide 3. The exploration is the kind of thing that
sits naturally in the \emph{exploration loop} sketched in
``Formalisation Second Steps'': a model is asked to investigate a
direction outward; if interesting, it produces a candidate
formalisation target; if formalisable, it becomes a Tide step. Here
the human played the picker role; in a more autonomous future the
exploration would itself be model-driven.

The exploration produced three figures (raw moments vs $t$, data
collapse onto $f_2(T)$, effective-rate sigmoid) which became inline
content in a new appendix in the primer
(\texttt{app:near-degenerate-crossover}, 46 pages with the new section
included).

\section{Time, cost, and skill iteration}

\paragraph{Wall-clock.} The three Tide steps and the side excursion
ran from morning to evening of 1 May 2026. Roughly:

\begin{itemize}
\item Tide 1 deliberation + formalisation + tide log + commit: $\sim 1.5$ hours.
\item Primer appendix for Tide 1 (with figures-free content): $\sim 0.5$ hours.
\item Tide 2 deliberation + formalisation + tide log + commit + primer
appendix subsection: $\sim 2.0$ hours.
\item Side excursion (script, figures, markdown note, primer appendix):
$\sim 1.0$ hour.
\item Tide 3 deliberation + formalisation + tide log + commit + primer
appendix update: $\sim 1.0$ hours.
\end{itemize}

Total: roughly six hours of attention, with a spread of $\sim 30$
minutes (Tide 3 formalisation) to $\sim 2$ hours (Tide 2 formalisation).

\paragraph{The skill itself.} The \texttt{tide} skill was authored at
the start of the afternoon, refined once mid-flight (after the
$L=x^2y^4$ flag, the integrability-or-not check was tightened). The
\texttt{lean-formalisation} skill was already mature; two small
additions to its ``Patterns and pitfalls'' section emerged during the
runs (above). No structural changes were needed.

\paragraph{Per-step costs.} Each Tide step cost $\$$10--$\$$30 in API
spend (estimated, from prior calibration), spread between Claude
session continuation (fairly cache-warm given the linear branch
structure), the deliberation GPT consult ($\sim \$2$ each), and one
GPT consult during Tide 1's formalisation (which timed out and so
charged the timeout cost). The aggregate was modest.

\paragraph{Skill compounding.} After three Tide steps, the
\texttt{tide} skill itself can be evaluated against actual usage:

\begin{itemize}
\item Step 1 (record direction): one-line, mechanical, never the
bottleneck.
\item Step 2 (deliberation): the load-bearing step. Caught real bugs in
$\sim 30\%$ of candidates. Worth the time.
\item Step 3 (formalisation hand-off): runs the
\texttt{lean-formalisation} skill, which absorbs the heavy lifting.
The hand-off note from Step 2 (file path, theorem names, proof
sketch) was used unchanged in all three runs.
\item Stop conditions: 4-round cap on deliberation never approached;
Step 3's ``runs to success'' condition was straightforward in all
three cases. The user's earlier instruction ``put off stop conditions
until they become a problem'' was vindicated.
\end{itemize}

\section{Tides 4 and 5: deepening and a cross-repo step}\label{sec:t4t5}

The original three excursions all formalised material in
\texttt{laplace}, building linearly on each other. Tides 4 and 5
extended the loop in two distinct directions.

\subsection*{Tide 4: rate-aware large-$T$ asymptotic for $f_1$}

Same evening as Tides 1--3 but after a brief detour through the
near-degenerate-crossover side excursion (\S\ref{sec:near-degenerate}),
which produced the universal scaling function $f_n(T) := \langle
y^{2n}\rangle^V_T$ for $V(y) = y^2/2 + y^4/24$ as a research artefact
sitting just outside the formal seabed. The note flagged a
``two-sided asymptotic for $f_n$'' as a candidate Tide step. After one
further Tide step (Tide 3, the universal scaling identity) had landed
the connection $\langle x^{2n}\rangle_{\mu_t}^{L^{\rm nd}_\varepsilon}
= \varepsilon^n f_n(\varepsilon^2 t)$, the natural next step was the
large-$T$ leading asymptotic of $f_1$:
\[
T \cdot f_1(T) \;=\; 1 + O(1/T) \qquad (T \to \infty).
\]

GPT recommended the rate-aware version (rather than a plain Tendsto)
because the seabed's $I_n$ asymptotics already carry rates; deriving
Tendsto from the rate is one line, deriving the rate from Tendsto is
not. GPT also flagged a strict-improvement candidate I had missed: an
all-$n$ proof via Gaussian rescaling $y = z/\sqrt T$, which bypasses
the seabed's $I_k$ ceiling at $k = 3$ entirely. We took the
seabed-bridge route as the smallest minimal step and recorded the
Gaussian-rescaling all-$n$ approach as a follow-up.

The proof bridges to the seabed's
\texttt{I\_0\_asymptotic} and \texttt{I\_2\_asymptotic} for the
anharmonic potential at the parameters $\alpha = 0$, $\lambda = \gamma
= 1$ (which is exactly $V$); the discriminant condition $\alpha^2 <
3\lambda\gamma$ becomes $0 < 3$, trivially. The bridge consists in
rescaling $D := \int e^{-TV}$ and $N := \int y^2 e^{-TV}$ to forms
$\sqrt T \cdot D$ and $T\sqrt T \cdot N$, each within $K/T$ of $a :=
\sqrt{2\pi}$, then bounding the ratio $(T\sqrt T \cdot N)/(\sqrt T
\cdot D)$ to $1$ by the standard ``small perturbation of constant''
estimate. 189 lines, no GPT consults during Step 3, single session.

The most useful incidental finding: Mathlib renamed
\texttt{div\_lt\_iff}, \texttt{div\_le\_iff},
\texttt{div\_le\_div\_iff} to their \texttt{$\cdot{}_0$} variants
(\texttt{div\_lt\_iff${}_0$} etc., adding explicit positivity
hypotheses); \texttt{push\_neg} is deprecated in favour of \texttt{push
Not}; \texttt{field\_simp} sometimes closes the goal outright and
sometimes leaves a polynomial residue requiring \texttt{ring} (the
diagnostic tells you which). Worth promoting to the
\texttt{lean-formalisation} skill's gotchas.

\subsection*{Tide 5: $\kappa_3$ vanishes under the harmonic Gibbs}

The next afternoon began with a survey of Tide-loop candidates across
the wider SRI portfolio (recorded as
\texttt{2026-05-02-tide-candidates-survey.md}), which produced five
options: the cumulant-vanishing in \texttt{threepoint}, a 2D
pure-quartic affine covariance in \texttt{laplace}, sextic 1D moments
also in \texttt{laplace}, the \texttt{sustm} factorisation
proposition, and a Gaussian/monomial precursor for the grammar paper.
The first won on Tide-shapedness: small, lands on an
existing-but-underused Lean repo, and demonstrates the loop on a
seabed other than \texttt{laplace}.

The headline statement:
\[
\kappa_3(\mu, L_{\rm harm}, x, t, x, x) = 0
\]
where $L_{\rm harm}(x) = (\lambda/2) x^2$ and the three observables in
the cumulant are all the linear function $x$. Mathematically, the
third cumulant of a centred Gaussian is zero; in the
\texttt{threepoint} formalism, this is one specific instance of the
abstract \texttt{kappa3} framework that had so far only been used in
its conditional form (every previous theorem in the repo was generic
over $(W, \mu, L, A)$, awaiting an instantiation).

GPT's most valuable contribution was negative: \emph{do not} attempt
to instantiate the \texttt{GibbsRegularity} or \texttt{GibbsObservable}
hypothesis bundles for the harmonic case. Those bundles encode the
differentiability needed for the FDT and cross-susceptibility
\emph{derivative} theorems. \texttt{kappa3} is defined at $h = 0$, so
only the value-at-zero matters. Unfolding \texttt{gibbsExp ... 0} and
simplifying $0 \cdot A w + L w = L w$ via \texttt{simp only [zero\_mul,
add\_zero]} reduces the proof to two parity integrals. Without that
advice, this Tide step would have grown into the kind of
``differentiation-under-the-integral plumbing'' detour that Tide
discipline is supposed to prevent. Estimated saving: $\sim 150$ lines.

The proof itself is light: one generic
\texttt{integral\_eq\_zero\_of\_odd} helper (using the standard $\int f
= -\int f$ argument from \texttt{integral\_neg\_eq\_self}), specialised
to the integrands $x \cdot e^{-tL}$ and $x^3 \cdot e^{-tL}$, then
\texttt{ring} to combine. 126 lines, no GPT consults during Step 3,
$\sim 1$ hour.

\subsection*{What was new about Tide 5}

The cross-repo step revealed nothing structurally surprising. The
\texttt{tide} skill, the \texttt{lean-formalisation} skill, the
deliberation-then-formalisation rhythm --- all of these transported
without modification. The seabed for \texttt{threepoint} is shaped
differently from \texttt{laplace}'s (a generic abstract framework
awaiting instantiation, rather than a concrete sequence of
asymptotics), and the GPT consult correctly diagnosed which parts of
that framework were load-bearing for our specific Tide step (none of
the analytic hypotheses) and which were not. The transport is real:
the loop is not coupled to any specific repo's conventions.

\section{Lessons}

Distilling the afternoon's experience into a small list:

\begin{enumerate}
\item \textbf{The deliberation step is where the value is.} GPT's
strict-improvement and bug-catching contributions during Step 2 saved
hours of Lean work in each run. The formalisation step is mechanical
when the deliberation has done its job.

\item \textbf{Linear branches beat parallel branches.} Branching each
Tide step's branch off the previous one's tip (\texttt{tide/quartic-moments}
$\to$ \texttt{tide/2d-semi-degenerate} $\to$ \texttt{tide/universal-scaling})
collapses the primer's commit pin to a single
\texttt{\textbackslash quarticCommit}. Parallel branches would have
required separate pins and a manual merge before the primer could
cite all three.

\item \textbf{Numerical sanity at the deliberation stage is
load-bearing.} Tide 1 caught the \emph{value} of the leading constant
$\sqrt{24}\,\Gamma(3/4)/\Gamma(1/4) \approx 1.6558$ to machine
precision before any Lean was written. Tide 2's data collapse and
Tide 3's substitution identity were each verified numerically
\emph{across multiple parameter values} before formalisation. Each
mismatch caught at this stage saves a wasted Lean session.

\item \textbf{Adopt GPT's strict-improvement candidate.} If GPT's
strengthened version of the candidate has the same proof, take it.
The reusable API matters more for follow-ups than the minimality of
any single step.

\item \textbf{Race short tactical questions, don't block on GPT.}
Tide 1's integrability question was sent to GPT and Claude in
parallel; Claude's local proof landed first. GPT's response timed out.
The lesson is in the \texttt{lean-formalisation} skill's ``When to
consult'' table: GPT for strategic and architectural questions, Claude
for tactical-Lean questions, and never block one on the other.

\item \textbf{Skill iteration is cheap; do it during the run.} The
\texttt{lean-formalisation} skill grew two small entries during the
afternoon, each landed in the run that produced them. Skills are not
finished products; they are the seabed, and they thicken with use.

\item \textbf{The exploration loop is the future.} The side excursion
between Tide 2 and Tide 3 produced a result (the universal scaling
collapse) which became Tide 3's direction. Wiring this kind of
exploration into the loop --- so that a model proposes the direction
without DM in the middle --- is the natural next ambition. Today the
human plays picker; tomorrow the picker is also a model.
\end{enumerate}

\section{Possible next steps}

\paragraph{All Tide branches now merged.} Tides 1--3 were merged to
\texttt{laplace/main} in three sequential PRs (\#2, replacement \#5,
\#4); Tide 4 was merged as PR \#6; Tide 5 was merged as PR \#1 in the
\texttt{threepoint} repo. The primer's
\texttt{\textbackslash quarticCommit} pin and
\texttt{\textbackslash leanrefQuartic} macro have been collapsed into
the standard \texttt{\textbackslash laplaceCommit}, which now points
at \texttt{44a6001} (the \texttt{laplace} main tip after Tide 4).

\paragraph{Continuation directions.} A candidates survey
(\texttt{2026-05-02-tide-candidates-survey.md}) inventories five
options drawn from across the SRI portfolio. The two next steps that
extend Tide 5's thread:
\begin{itemize}
\item Abstract parity ($\kappa_3$ vanishes for any odd-observable
triple under negation-invariant $\mu$ and even $L$). Strict
generalisation of Tide 5; ${\sim}50$ lines.
\item Four-point cumulant ($\kappa_4$) for the harmonic Gibbs.
Pairs naturally with $\kappa_3$ and the flow equation; closes the
cumulant story for the Gaussian case.
\end{itemize}
Other directions still open:
\begin{itemize}
\item Two-sided asymptotic for $f_n(T)$ in the universal potential at
small $T$ (Candidate~B from the Tide~4 deliberation); requires a
uniform Taylor expansion against the quartic Gibbs.
\item All-$n$ large-$T$ asymptotic via Gaussian rescaling rather than
seabed $I_n$ extension (the strict-improvement candidate from the
Tide~4 deliberation; the natural follow-up to Tide~4).
\item 2D pure-quartic affine covariance and sextic 1D moments
(candidates~2 and~3 from the survey); both are mechanical extensions.
\item \texttt{sustm} Proposition~1 (factorisation for mismatched
input--component pairs); requires a fresh seabed for noisy Turing
machine susceptibilities.
\item Grammar-paper precursors (Gaussian or single-monomial $Z_n[\phi]$
expansion); the start of a planned multi-Tide programme rather than a
single excursion.
\end{itemize}

\paragraph{Promote pieces to \texttt{lean-common}.} The 2D Gibbs
wrappers from Tide 2 (\texttt{partitionFunction2D},
\texttt{gibbsExpectation2D}, \texttt{gibbsCov2D}) are promotion
candidates as soon as a second project needs them. Likewise
\texttt{quartic\_integrable\_pow} (Tide 1) generalises trivially to
arbitrary $L = (\varepsilon/2)x^2 + (\gamma/24) x^4$ once the second
use materialises.

\paragraph{Wire the exploration loop.} A model-driven
\emph{exploration} step --- ``starting from this seabed, propose three
candidate directions outward and check each numerically'' --- followed
automatically by a Tide step on the most promising one would close the
loop without DM in the middle. The numerical-check infrastructure used
in the side excursion (a few-line scipy script) is the existence
proof that this is a small step from where the loop is now.

\paragraph{Adopt \texttt{leanblueprint}.} The earlier retrospective
already argued for this once a second result was formalised. By the
end of this two-afternoon stretch there are six formalised primer
claims (three from the morning of 27 April plus its multi-day
follow-up, plus four more from days 1--2 of May: the 1D quartic affine
covariance, the 2D semi-degenerate affine covariance, the universal
scaling identity, and the rate-aware large-$T$ asymptotic of the
universal $f_1$\footnote{Numbered
\texttt{lem:quartic-cov-affine}, \texttt{lem:semi-degenerate-cov-affine},
\texttt{lem:nd-scaling}, and the inline-referenced
\texttt{gibbsExpectation\_universal\_sq\_large\_t\_rate} in the
\texttt{near-degenerate-crossover} appendix.}), with margin markers
across two repos (\texttt{laplace} and \texttt{threepoint}).
\texttt{leanblueprint}'s dependency-graph view would now be genuinely
useful, and the cross-repo coordination would benefit from the
shared-blueprint workflow.

\section{Acknowledgements}

\begin{itemize}
\item GPT-5.5~Pro provided five deliberation consultations preserved
verbatim under the relevant project's \texttt{tide-log/} directory (\texttt{sri/projects/primer/tide-log/} for tides 1--4, \texttt{sri/projects/patterning/tide-log/} for tide 5):
\texttt{gpt55\_tide\_quartic\_v1.md},
\texttt{gpt55\_tide\_2d\_v1.md},
\texttt{gpt55\_tide\_scaling\_v1.md},
\texttt{gpt55\_tide\_fn\_v1.md}, and
\texttt{gpt55\_tide\_kappa3\_v1.md}.
\item The \texttt{tide} skill is at
\texttt{sri/.claude/skills/tide/SKILL.md}; the
\texttt{lean-formalisation} skill is at
\texttt{sri/.claude/skills/lean-formalisation/SKILL.md}. Both are
short and read-once; they are deliberately thin.
\item Geoffrey Irving's \href{https://github.com/girving/aks}{aks}
repository continues to be the model for project layout, audit
discipline, and \texttt{CLAUDE.md} structure (transitively, via the
seabed).
\item The mathematical content sits in the SLT Susceptibility Primer
(Elliott--Murfet, 2026), which now carries three appendices on the
degenerate-case excursions plus updated provenance paragraphs
crediting the Tide loop.
\item The formalisation depends throughout on
\href{https://github.com/leanprover-community/mathlib4}{Mathlib}.
\end{itemize}

\end{document}
